\section{Exact Parry Divergence and Sharp Support Barriers}
\label{sec:parry-divergence}

The Parry measure is the unique maximal-entropy equilibrium state of the
ambient golden-mean shift, whose topological entropy is
\(h_{\mathrm{top}}(X)=\log\varphi\). This section quantifies how far
folded rotation laws sit from ambient equilibrium when their effective support
is much smaller than the full shift language. The exact identities below rely
on endpoint collapse for the golden-mean Parry cylinders, and no comparable
closed formula is claimed for arbitrary shifts of finite type.

\begin{definition}[R\'enyi entropy and divergence]\label{def:renyi-entropy-divergence}
Let \(E\) be a finite set. For a probability law \(\nu\) on \(E\), define
\[
H_q(\nu):=
\begin{cases}
\log |\supp(\nu)|,
& q=0,\\[0.4em]
\dfrac{1}{1-q}\log\!\Bigl(\sum_{x\in E}\nu(x)^q\Bigr),
& q\in(0,\infty)\setminus\{1\},\\[0.8em]
-\sum_{x\in E}\nu(x)\log \nu(x),
& q=1,\\[0.4em]
-\log \max_{x\in E}\nu(x),
& q=\infty.
\end{cases}
\]
For probability laws \(\nu\) and \(\rho\) on \(E\), define
\[
D_q(\nu\|\rho):=
\begin{cases}
-\log \rho(\supp(\nu)),
& q=0,\\[0.4em]
\dfrac{1}{q-1}\log\!\Bigl(\sum_{x\in E}\nu(x)^q\rho(x)^{1-q}\Bigr),
& q\in(0,\infty)\setminus\{1\},\\[0.8em]
\KL(\nu\|\rho),
& q=1,\\[0.4em]
\log \max_{x:\,\nu(x)>0}\dfrac{\nu(x)}{\rho(x)},
& q=\infty,
\end{cases}
\]
with the convention \(0\log 0:=0\), and with \(D_q(\nu\|\rho):=\infty\) for
\(q>0\) if \(\nu\) is not absolutely continuous with respect to \(\rho\).
\end{definition}

\subsection{Exact Parry Identity}

\begin{theorem}[Exact divergence identity relative to the Parry measure]\label{thm:parry-kl-identity}
Let \(\nu\) be a probability measure on \(X_m\), and define the endpoint masses
\[
\eta_{ij}(\nu):=\nu\bigl(\{u\in X_m:u_1=i,\ u_m=j\}\bigr),
\qquad i,j\in\{0,1\}.
\]
If
\[
H(\nu):=-\sum_{u\in X_m}\nu(u)\log \nu(u)
\]
denotes the Shannon entropy of \(\nu\), then
\[
\KL(\nu\|\pi_m)
=
m\log\varphi-H(\nu)-\log\pi(0)
{}+(\eta_{11}(\nu)-\eta_{00}(\nu))\log\varphi.
\]
Consequently,
\[
m\log\varphi-H(\nu)-\log\pi(0)-\log\varphi
\le
\KL(\nu\|\pi_m)
\le
m\log\varphi-H(\nu)-\log\pi(0)+\log\varphi.
\]
\end{theorem}

\begin{proof}
By Proposition~\ref{prop:parry-endpoint-collapse},
\[
\log \pi_m(u)
=
\log\pi(0)-m\log\varphi
{}+\bigl(\mathbf{1}_{\{u_1=0\}}-\mathbf{1}_{\{u_m=1\}}\bigr)\log\varphi.
\]
Therefore
\[
\begin{aligned}
\KL(\nu\|\pi_m)
&=
\sum_{u\in X_m}\nu(u)\log\frac{\nu(u)}{\pi_m(u)}\\
&=
-H(\nu)-\sum_{u\in X_m}\nu(u)\log\pi_m(u)\\
&=
m\log\varphi-H(\nu)-\log\pi(0)
{}+\bigl(\eta_{11}(\nu)-\eta_{00}(\nu)\bigr)\log\varphi.
\end{aligned}
\]
Since \(\eta_{11}(\nu)-\eta_{00}(\nu)\in[-1,1]\), the two-sided bound follows.
\end{proof}

\begin{remark}[Specificity of the endpoint collapse]\label{rem:endpoint-collapse-specificity}
Theorem~\ref{thm:parry-kl-identity} depends essentially on
Proposition~\ref{prop:parry-endpoint-collapse}. For a general shift of finite
type, Parry cylinder weights are controlled by boundary states in an edge
presentation but need not collapse to a scalar endpoint statistic of the kind
available here. The exact identity is part of what makes the golden-mean
setting especially rigid.
\end{remark}

\begin{corollary}[No fold-entropy loss on the Sturmian slice]\label{cor:no-fold-entropy-loss}
Assume \(\beta=\alpha\). Then
\[
\KL(\pi_m^{\mathrm{fold}}\|\pi_m)
=
m\log\varphi-H(\mu_m)-\log\pi(0)
{}+\bigl(\eta_{11}(\pi_m^{\mathrm{fold}})
-\eta_{00}(\pi_m^{\mathrm{fold}})\bigr)\log\varphi.
\]
Likewise, for every \(N\ge 1\),
\[
\KL(\hat\pi_{m,N}\|\pi_m)
=
m\log\varphi-H(\hat\mu_{m,N})-\log\pi(0)
{}+\bigl(\eta_{11}(\hat\pi_{m,N})
-\eta_{00}(\hat\pi_{m,N})\bigr)\log\varphi.
\]
\end{corollary}

\begin{proof}
Apply Theorem~\ref{thm:parry-kl-identity} with
\(\nu=\pi_m^{\mathrm{fold}}\) and \(\nu=\hat\pi_{m,N}\), and use
Corollary~\ref{cor:sturmian-exact-relabeling} to replace
\(H(\pi_m^{\mathrm{fold}})\) by \(H(\mu_m)\) and
\(H(\hat\pi_{m,N})\) by \(H(\hat\mu_{m,N})\).
\end{proof}

\begin{proposition}[Order-uniform R\'enyi window relative to the Parry measure]
\label{prop:parry-renyi-window}
For every probability law \(\nu\) on \(X_m\) and every \(q\in[0,\infty]\),
\[
m\log\varphi-H_q(\nu)-\log\pi(0)-\log\varphi
\le
D_q(\nu\|\pi_m)
\le
m\log\varphi-H_q(\nu)-\log\pi(0)+\log\varphi.
\]
\end{proposition}

\begin{proof}
For \(q=1\), this is Theorem~\ref{thm:parry-kl-identity}. For \(q=0\), let
\(S:=\supp(\nu)\). Proposition~\ref{prop:parry-endpoint-collapse} gives
\[
\pi(0)\varphi^{-(m+1)}
\le
\pi_m(u)
\le
\pi(0)\varphi^{-(m-1)}
\qquad (u\in X_m),
\]
and therefore
\[
|S|\pi(0)\varphi^{-(m+1)}
\le
\pi_m(S)
\le
|S|\pi(0)\varphi^{-(m-1)}.
\]
Since \(D_0(\nu\|\pi_m)=-\log \pi_m(S)\) and \(H_0(\nu)=\log|S|\), the
displayed bounds follow.

For \(q=\infty\), the same pointwise bounds on \(\pi_m(u)\) imply
\[
\frac{\varphi^{m-1}}{\pi(0)}\max_{u\in X_m}\nu(u)
\le
\max_{u:\,\nu(u)>0}\frac{\nu(u)}{\pi_m(u)}
\le
\frac{\varphi^{m+1}}{\pi(0)}\max_{u\in X_m}\nu(u).
\]
Taking logarithms and using \(H_\infty(\nu)=-\log\max_u \nu(u)\) gives the
claim.

Now fix \(q\in(0,\infty)\setminus\{1\}\). By
Proposition~\ref{prop:parry-endpoint-collapse},
\[
\pi_m(u)
=
\pi(0)\varphi^{-m+\theta(u)},
\qquad
\theta(u):=\mathbf{1}_{\{u_1=0\}}-\mathbf{1}_{\{u_m=1\}}\in[-1,1].
\]
Hence
\[
\sum_{u\in X_m}\nu(u)^q\pi_m(u)^{1-q}
=
\pi(0)^{1-q}\varphi^{m(q-1)}
\sum_{u\in X_m}\nu(u)^q\varphi^{(1-q)\theta(u)}.
\]
Because \(\theta(u)\in[-1,1]\), one has
\[
\varphi^{-|1-q|}\sum_{u\in X_m}\nu(u)^q
\le
\sum_{u\in X_m}\nu(u)^q\varphi^{(1-q)\theta(u)}
\le
\varphi^{|1-q|}\sum_{u\in X_m}\nu(u)^q.
\]
Taking logarithms, dividing by \(q-1\), and using the definition of \(H_q\)
gives the same two-sided bound with error at most \(\log\varphi\).
\end{proof}

\subsection{Support Conditioning and Endpoint Class Counts}

\begin{proposition}[Support conditioning for R\'enyi divergence]
\label{prop:renyi-support-conditioning}
Let \(E\) be finite, let \(\rho\) be a probability law on \(E\), and let
\(S\subset \supp(\rho)\) satisfy \(\rho(S)>0\). Writing
\(\rho^S:=\rho(\cdot\mid S)\), every probability law \(\nu\) on \(E\) with
\(\supp(\nu)\subset S\) satisfies
\[
D_q(\nu\|\rho)
=
D_q(\nu\|\rho^S)-\log \rho(S)
\qquad (q\in[0,\infty]).
\]
Consequently,
\[
\inf_{\substack{\nu\in\Delta(E)\\ \supp(\nu)\subset S}} D_q(\nu\|\rho)
=
-\log \rho(S)
\qquad (q\in[0,\infty]).
\]
If \(q\in(0,\infty]\), the minimizer is unique and equals \(\rho^S\). If
\(q=0\), the minimizers are exactly the laws whose support is \(S\).
\end{proposition}

\begin{proof}
For \(q=1\), the identity is the usual conditioning formula for relative
entropy. For \(q=0\),
\[
D_0(\nu\|\rho)
=
-\log \rho(\supp(\nu))
=
-\log \rho^S(\supp(\nu))-\log \rho(S)
=
D_0(\nu\|\rho^S)-\log \rho(S).
\]
For \(q=\infty\),
\[
D_\infty(\nu\|\rho)
=
\log\max_{x:\,\nu(x)>0}\frac{\nu(x)}{\rho(x)}
=
\log\max_{x:\,\nu(x)>0}\frac{\nu(x)}{\rho(S)\rho^S(x)}
=
D_\infty(\nu\|\rho^S)-\log \rho(S).
\]
Finally, for \(q\in(0,\infty)\setminus\{1\}\),
\[
\sum_{x\in E}\nu(x)^q\rho(x)^{1-q}
=
\sum_{x\in S}\nu(x)^q\bigl(\rho(S)\rho^S(x)\bigr)^{1-q}
=
\rho(S)^{1-q}\sum_{x\in S}\nu(x)^q(\rho^S(x))^{1-q},
\]
which gives the claimed identity after taking logarithms and dividing by
\(q-1\).

For every \(q>0\), the standard nonnegativity \(D_q(\nu\|\rho^S)\ge 0\) shows
that \(D_q(\nu\|\rho)\ge -\log\rho(S)\). Equality holds exactly when
\(D_q(\nu\|\rho^S)=0\), namely when \(\nu=\rho^S\). For \(q=0\),
\[
D_0(\nu\|\rho^S)=-\log \rho^S(\supp(\nu))\ge 0,
\]
with equality if and only if \(\rho^S(\supp(\nu))=1\), equivalently
\(\supp(\nu)=S\).
\end{proof}

\begin{lemma}[Endpoint class counts]\label{lem:endpoint-class-counts}
Extend the Fibonacci sequence by \(F_0:=0\), \(F_1:=1\), and
\(F_{n+2}=F_{n+1}+F_n\). For \(m\ge 2\), let
\[
X_m^{ij}:=\{u\in X_m:u_1=i,\ u_m=j\},
\qquad i,j\in\{0,1\}.
\]
Then
\[
|X_m^{00}|=F_m,\qquad
|X_m^{01}|=|X_m^{10}|=F_{m-1},\qquad
|X_m^{11}|=F_{m-2}.
\]
\end{lemma}

\begin{proof}
The admissible length-\(m\) words in \(X_m\) are the paths of length \(m-1\) in
the golden-mean graph with adjacency matrix
\[
A=\begin{pmatrix}1&1\\1&0\end{pmatrix}.
\]
For the standard Fibonacci extension above,
\[
A^n=
\begin{pmatrix}
F_{n+1} & F_n\\
F_n & F_{n-1}
\end{pmatrix}
\qquad (n\ge 1).
\]
Hence the number of admissible words of length \(m\) with endpoints \((i,j)\) is
\((A^{m-1})_{ij}\), which gives the displayed formulas.
\end{proof}

\subsection{Sharp Support Barriers}

\begin{theorem}[Sharp support barrier against the Parry measure]\label{thm:parry-support-barrier}
Let \(m\ge 2\) and \(1\le k\le F_{m+2}\). Define
\[
M_m(k):=\max\{\pi_m(S):S\subset X_m,\ |S|\le k\}.
\]
Then
\[
M_m(k)=
\begin{cases}
k\,\pi(0)\varphi^{-(m-1)},
& 1\le k\le F_m,\\[0.4em]
F_m\pi(0)\varphi^{-(m-1)}+(k-F_m)\pi(0)\varphi^{-m},
& F_m<k\le F_m+2F_{m-1},\\[0.4em]
1-(F_{m+2}-k)\pi(0)\varphi^{-(m+1)},
& F_m+2F_{m-1}<k\le F_{m+2}.
\end{cases}
\]
Moreover,
\[
\inf_{\substack{\nu\in\Delta(X_m)\\ |\supp(\nu)|\le k}}D_q(\nu\|\pi_m)
=
-\log M_m(k),
\qquad q\in[0,\infty],
\]
and
\[
\inf_{\substack{\nu\in\Delta(X_m)\\ |\supp(\nu)|\le k}}\TV(\nu,\pi_m)
=
1-M_m(k).
\]
If \(q\in(0,\infty]\), the R\'enyi infimum is attained uniquely by conditioning
\(\pi_m\) on a maximizing set \(S\subset X_m\); equivalently, one takes all
words from the heaviest endpoint classes until the budget \(k\) is exhausted.
If \(q=0\), the minimizers are exactly the laws whose support is such a
maximizing set. The same conditioned Parry laws attain the TV infimum.
\end{theorem}

\begin{proof}
The argument has three layers. One first orders individual Parry cylinder
weights by endpoint class, then maximizes the mass of a \(k\)-point set by a
greedy filling of these endpoint classes, and finally invokes
Proposition~\ref{prop:renyi-support-conditioning} to identify the optimal
R\'enyi law on a fixed support.

By Proposition~\ref{prop:parry-endpoint-collapse}, every word in \(X_m^{00}\) has
Parry weight \(\pi(0)\varphi^{-(m-1)}\), every word in \(X_m^{01}\cup X_m^{10}\)
has weight \(\pi(0)\varphi^{-m}\), and every word in \(X_m^{11}\) has weight
\(\pi(0)\varphi^{-(m+1)}\). These three values are strictly ordered. Together
with Lemma~\ref{lem:endpoint-class-counts}, this implies that the largest
possible Parry mass of a \(k\)-point subset is obtained greedily by exhausting
the endpoint classes in the order
\[
X_m^{00},\qquad X_m^{01}\cup X_m^{10},\qquad X_m^{11},
\]
which gives the formula for \(M_m(k)\).

Now let \(S\subset X_m\) and let \(\nu\) be supported on \(S\). By
Proposition~\ref{prop:renyi-support-conditioning},
\[
D_q(\nu\|\pi_m)
\ge
-\log \pi_m(S)
\qquad (q\in[0,\infty]),
\]
with equality for the conditioned law \(\pi_m^S:=\pi_m(\cdot\mid S)\), and
with uniqueness when \(q\in(0,\infty]\). Therefore
\[
\inf_{\substack{\nu\in\Delta(X_m)\\ |\supp(\nu)|\le k}}D_q(\nu\|\pi_m)
=
-\log\max_{\substack{S\subset X_m\\ |S|\le k}}\pi_m(S)
=
-\log M_m(k)
\]
for every \(q\in[0,\infty]\). The characterization of minimizers follows from
the same proposition.

For total variation, if \(\nu\) is supported on \(S\), then
\[
\TV(\nu,\pi_m)
\ge
\pi_m(S^c)
=
1-\pi_m(S).
\]
Equality holds for \(\nu=\pi_m^S\), since then
\[
\TV(\pi_m^S,\pi_m)
=
\frac12\sum_{u\in X_m}\bigl|\pi_m^S(u)-\pi_m(u)\bigr|
=
1-\pi_m(S).
\]
Taking the best support set \(S\) gives the total-variation formula.
\end{proof}

\begin{corollary}[Exact complexity-divergence law]\label{cor:exact-complexity-divergence-law}
Let \(m\ge 2\) and \(1\le k\le F_m\). Then
\[
\inf_{\substack{\nu\in\Delta(X_m)\\ |\supp(\nu)|\le k}}D_q(\nu\|\pi_m)
=
(m-1)\log\varphi-\log\pi(0)-\log k,
\qquad q\in[0,\infty],
\]
and
\[
\inf_{\substack{\nu\in\Delta(X_m)\\ |\supp(\nu)|\le k}}\TV(\nu,\pi_m)
=
1-k\pi(0)\varphi^{-(m-1)}.
\]
If \(q\in(0,\infty]\), the R\'enyi minimizers are exactly the uniform laws on
\(k\)-point subsets of \(X_m^{00}\). If \(q=0\), the minimizers are exactly the
laws whose support is such a \(k\)-point subset. The TV minimizers are the
conditioned Parry laws on these same subsets.

Consequently, if \(k_m\le F_m\) and
\[
\lim_{m\to\infty}\frac{1}{m}\log k_m=\gamma\in[0,\log\varphi],
\]
then
\[
\lim_{m\to\infty}
\inf_{\substack{\nu\in\Delta(X_m)\\ |\supp(\nu)|\le k_m}}
\frac{1}{m}D_q(\nu\|\pi_m)
=
\log\varphi-\gamma
=
h_{\mathrm{top}}(X)-\gamma,
\qquad q\in[0,\infty].
\]
In particular, every subexponential support budget still incurs the full rate
\(\log\varphi\), and cancelling the Parry penalty requires support growth of
order \(\varphi^m\).
\end{corollary}

\begin{proof}
The first branch of Theorem~\ref{thm:parry-support-barrier} gives the displayed
formulas. The description of the minimizers follows because all words in
\(X_m^{00}\) carry the same Parry weight \(\pi(0)\varphi^{-(m-1)}\), so the
conditioned Parry law on any \(k\)-point subset of \(X_m^{00}\) is uniform.
The limit statement is obtained by dividing the R\'enyi formula by \(m\) and
letting \(m\to\infty\).
\end{proof}
